\documentstyle[%


righttag%


,11pt%


,amssymb%


,russian%               remove in Eng.


,izvru%                 Set izven% in Eng.


]{amsart}





% NB: big piece after EOD





\newcommand{\rank}{\operatorname{rank}}


\newcommand{\tts}[1]{}


\renewcommand{\baselinestretch}{0.97}





%\hoffset=-15mm%                  For Eng.


\setcounter{page}{11}


\god{2002}


\nomer{2~(477)} %                        Remove in Eng.


%\nomer{Vol.\,46, No.\,1}%               Set in Eng.


%\str{\pageref{begin}--\pageref{end}}%  Set in Eng.


\udk{517.977}





\author{\it Л.Т.\,Ащепков, Д.В.\,Давыдов}


% NB:   ^^ remove in Eng.


\title{Стабилизация наблюдаемой линейной системы управления


с постоянными интервальными коэффициентами}





\date{}


%\pagestyle{myheadings}%         Set in Eng.


\pagestyle{plain} %              Remove in Eng.


\begin{document}


%\thispagestyle{plain}%          Set in Eng.


\maketitle


\label{begin}








Методами упреждающего управления (predictive control) и интервального


анализа разработана процедура стабилизации системы линейных автономных


дифференциальных уравнений с интервальными коэффициентами, включающая этапы


пассивного наблюдения фазовых состояний и активного управления.


Приведены достаточные условия притяжения всех траекторий системы к положению


равновесия.





\section*{Введение}





Задача стабилизации линейной системы управления по наблюдениям фазовых


состояний в детерминированной и стохастической постановке относится к ряду


классических [1]. Новые интересные аспекты в ее трактовке и решении, не


укладывающиеся в рамки традиционной теории, возникают при интервальной


неопределенности коэффициентов в уравнениях движения и наблюдения [2].


Источниками интервальности могут быть неполнота знаний об объекте


управления и вытекающие отсюда ошибки моделирования, погрешности


вычисления коэффициентов или последствия линеаризации нелинейных уравнений с


неопределенными параметрами.





Потенциальный континуум реализаций неопределенных коэффициентов в интервалах


задания составляет основное отличие интервальной задачи от детерминированной.


Вместе с тем, полное отсутствие информации о реализациях коэффициентов не


позволяет трактовать ее как стохастическую и использовать хорошо разработанную


теорию фильтрации [1].





Применяемый здесь подход к решению интервальной задачи стабилизации близок


статье [3]. Основная движущая идея состоит в построении детерминированного


программно-позиционного управления, которое на конечных отрезках времени


действует как оператор сжатия в фазовом пространстве, ``сужая'' интегральную


воронку траекторий каждой системы из рассматриваемого континуума систем.


Если циклически чередовать по времени восстановление фазовых векторов с


управлением, то при определенных условиях удается обеспечить асимптотическую


устойчивость положения равновесия системы в целом (в большом) независимо


от реализаций ее коэффициентов в заданных интервалах.





Процедура стабилизации ориентирована на реальное время.


Основные трудоемкие операции --- это вычисление матричной экспоненты и


интеграла от линейно преобразованных результатов наблюдения фазовой


траектории на конечных отрезках времени. Вспомогательные операции обращения


двух постоянных матриц могут быть выполнены заблаговременно.





Процедура и достаточное условие стабилизации интервальной системы являются


основными результатами данной работы. Достаточное условие формулируется


через известные середины и длины интервалов неопределенности и в этом


смысле конструктивно. Оно заведомо выполняется для интервалов малой длины,


если эталонная ``центральная'' система (коэффициенты которой соответствуют


серединам интервалов) обладает свойствами управляемости и наблюдаемости. С


точки зрения центральной системы достаточное условие стабилизируемости


можно трактовать как критерий ее робастности по отношению к применяемой


процедуре стабилизации.





\section{Постановка задачи}





Рассмотрим наблюдаемую систему управления


\begin{gather}


\Dot x = Ax+Bu, \\


y(t)=Cx(t), \ \ t \ge 0,


\end{gather}


где


$x$, $u$, $y$ --- векторы фазового состояния, управления и наблюдения из


$R^n$, $R^r$, $R^m$; $A$, $B$, $C$ ---


неопределенные постоянные матричные коэффициенты размерности


$n \times n$, $n \times r$, $m \times n$, $m \le n$,


со значениями из замкнутых интервалов


\begin{equation}


|A-A_0| \le \Delta A, \ \ |B-B_0| \le \Delta B, \ \ |C-C_0| \le \Delta C;


\end{equation}


$A_0$, $B_0$, $C_0$ и $\Delta A$, $\Delta B$, $\Delta C$


--- матрицы соответствующих размерностей с заданными произвольными и


неотрицательными элементами. Модули матриц и матричные неравенства в


(3) и далее понимаются поэлементно.





Матрицы $A_0$, $B_0$ предполагаются управляемыми, а матрицы $A_0$, $C_0$


--- наблюдаемыми. Это означает [4] выполнение алгебраических условий


\begin{gather}


\rank (B_0, A_0B_0, \dotsc, A^{n-1}_0B_0) = n,\\


\rank (C_0', A_0'C_0', \dotsc, ({A'}_0)^{n-1} C_0') = n


\end{gather} 


на ранги блочных матриц. Символ $'$ (штрих) везде означает операцию


транспонирования.





Требуется: 1) найти кусочно-непрерывное по $t$


на конечных отрезках времени программно-позиционное управление $u(x_0, t)$,


обеспечивающее притяжение траекторий замкнутой системы


\begin{equation}


\Dot x = Ax+Bu(x_0,t), \ \ x(0)=x_0


\end{equation}%6


к положению равновесия $x=0$ для любых начальных векторов $x_0$ из $R^n$ и


любых матриц $A$, $B$ из интервалов (3); 2) организовать процедуру построения


управления $\wh u(x_0, t) \approx u(x_0, t)$ с такими же свойствами, как у


$ u(x_0, t)$, по результатам наблюдения $ y(t)$ траектории $ x(t) $ системы


\begin{equation}


\Dot x = Ax+B \wh u(x_0,t), \ \ y(t)=Cx(t), \ \ t \ge 0;


\end{equation}


3) сформулировать достаточные условия стабилизируемости системы (7)


независимо от выбора начального вектора $x_0$ из $R^n$ и матриц


$A$, $B$, $C$ из интервалов (3).





\section{Стабилизирующее управление}





Разберем сначала относительно простой случай, когда все координаты фазового


вектора доступны непосредственному измерению в моменты времени


$t=0,T,2T,\dotsc$ с постоянным шагом $T >0$. Для упрощения можно считать


матрицу $C=E$ в уравнении (2) единичной.





Систему (1), (2) при $A=A_0$, $B=B_0$, $C=C_0$ назовем центральной.





Потребуем, чтобы стабилизирующее управление обладало следующими двумя


свойствами:


а) переводило траекторию центральной системы из любого состояния


$x(0)=x_0$ в состояние $x(T)=0$ за время $ T>0$; б) имело минимальную норму


$$


\| u \| = \bigg(\int_{0}^{T} u(t)' u(t) \, dt \bigg)^{1/2}


$$


в пространстве $L^r_2(0,T)$ векторных функций размерности $r$


с суммируемым на отрезке $[0,T]$ квадратом евклидовой нормы.





Из теории линейных систем известно [5], что в предположении (4) существует


единственное управление


\begin{equation}


u(t)=B_0'e^{(T-t)A_0'} z, \ \ 0 \le t \le T,


\end{equation} 


с требуемыми свойствами. Здесь $e^{(T-t)A_0'}$ --- матричная экспонента [6],


$z$ --- решение системы линейных алгебраических уравнений


\begin{equation}


D_0 z = f_0


\end{equation}


с матричными коэффициентами


\begin{equation}


D_0 = \int_0^T e^{\tau A_0}B_0 B_0'e^{\tau A_0'} \, d \tau,


\quad f_0 = - e^{T A_0}x_0


\end{equation}


размерности $n \times n$ и $n \times 1$ соответственно. Условие (4)


гарантирует ([1], c.\,49) невырожденность матрицы $D_0$.





Опираясь на результаты работы [8], опишем действие управления (8)--(10) на


состояние $x(T)$ системы (1) с начальным условием $x(0)=x_0$ при


неопределенных матрицах $A$, $B$. По формуле Коши ([9], c.\,134)


с использованием обозначений


\begin{equation}


D = \int_0^T e^{\tau A}B B_0'e^{\tau A_0'} \, d \tau,


\quad f = - e^{T A}x_0


\end{equation}


имеем


$$


x(T)= e^{TA}x_0+\int_0^T e^{(T-\tau)A}B u(t) \, dt= -f+Dz.


$$


Отсюда и из (9) следует


\begin{equation}


|x(T)|=|Dz-f-(D_0z-f_0)| \le |D-D_0|\,|z|+|f-f_0| \le \Delta D |z|+\Delta f,


\end{equation}


где


\begin{gather}


\Delta D = \int_{0}^{T}


[(e^{ \tau ( |A_0|+\Delta A ) } -e^{\tau |A_0|})\,


( |B_0|+\Delta B )+| e^{ \tau A_0 } | \Delta B]


\, |B_0' e^{ \tau A_0' }| \, d\tau, \\


\Delta f =


( e^{ T(|A_0|+\Delta A)} - e^{T|A_0|}) |x_0|.


\end{gather} 





Из (9), (10) имеем


$$


|z|= |D_0^{-1}f_0| \le |D_0^{-1}e^{TA_0}|\,|x_0|,


$$


поэтому верхняя оценка (12) с учетом (14) примет вид


\begin{equation}


|x(T)| \le M |x_0|, \quad


M=\Delta D |D_0^{-1}e^{TA_0}|+e^{ T(|A_0|+\Delta A)} - e^{T|A_0|}.


\end{equation} 





Очевидно, если в неравенстве (15) норма матрицы


$M$ меньше единицы, то для согласованных с нею норм векторов получим условие


$\|x(T)\|< \|x_0\|$. Это означает, что при $\|M \| < 1$


управление (8)--(10) действует как оператор сжатия в фазовом пространстве,


переводя начальные точки $x_0$ системы (1) с неопределенными коэффициентами


$A$, $B$ за время $T$ в более близкие к началу координат точки $x(T)$.





Mногократное повторение операции ``сжатия'' фазового пространства


за счет последовательного применения управлений типа (8)--(10) на отрезках


времени $[T,2T], [2T,3T],\dotsc$ вызовет притяжение траектории системы (1)


к началу координат, т.\,е. $x(t) \rightarrow 0$ при $t \rightarrow \infty$.


Действительно, если принять


\begin{gather}


u(t)=B_0'e^{(kT-t)A_0'} z^{k-1}, \ \ (k-1)T \le t \le kT, \\


%


D_0 z^{k-1} = - e^{TA_0} x((k-1)T), \ \  k=1,2,\dotsc,\notag


\end{gather}


то для соответствующего решения $x(t)$


системы (6) легко устанавливается аналог неравенства (15)


$$


|x(kT)| \le M |x((k-1)T)|, \ \ k=1,2, \dotsc,


$$


с той же матрицей $M$ на каждом из отрезков $[(k-1)T, kT]$.





Следовательно, неравенство $\|M \| < 1$


для матрицы (15) вместе с предположением (4) гарантирует стабилизацию


траекторий системы (6) независимо от выбора начального вектора


$x_0$ из $R^n$ и матриц $A$, $B$ из интервалов (3).





\section{Оценка фазового состояния}





Закон управления (16) предполагает известными состояния


$x((k-1)T)$, $k=1,2,\dotsc$,


системы (6) и, вообще говоря, не реализуем при мультипликативных


неопределенностях в уравнениях наблюдения (2). Тем не менее, при определенном


согласовании параметров $T$, $\Delta A$, $\Delta B$, $\Delta C$


с точностью восстановления векторов $x((k-1)T)$, $k=1,2,\dotsc$,


управление (16) в модифицированной форме сохраняет стабилизирующее свойство.


Перейдем к точным формулировкам.





Опишем двухфазную процедуру управления на основном отрезке времени


$[0,T]$. На первой фазе $(0 \le t \le \theta) $


процедуры по свободному движению системы (1), (2) при $u \equiv 0$


строится оценка $\wh x_0$ ее неизвестного начального состояния $x_0=x(0)$.


На второй фазе $(\theta \le t \le T)$ производится активное приближение


состояния $x(T)$ системы (1) к началу координат управлением типа (8) с


использованием найденной оценки $\wh x_0$.





На первой фазе, полагая в уравнениях (1), (2) $ u \equiv 0$,


приходим к однородной системе наблюдения


\begin{equation}


\Dot x = Ax, \ \ y(t)=Cx(t), \ \  0 \le t \le \theta.


\end{equation}


Отсюда с использованием формулы Коши известные измерения $y(t)$


можно выразить через неизвестные параметры $A$, $C$, $x_0$


\begin{equation}


y(t)= C e^{tA} x_0, \ \ 0 \le t \le \theta.


\end{equation}





Применим к измерениям (18) линейное интегральное преобразование с матрицей


$e^{tA_0'} C_0'$. Тогда получим систему линейных алгебраических уравнений


\begin{equation}


W x_0 = g


\end{equation}


относительно вектора $x_0$ с матричными коэффициентами


\begin{equation}


W = \int_0^{\theta} e^{t A_0'} C_0' C e^{ t A} \, d t,\quad


g = \int_0^{\theta} e^{t A_0'} C_0' y(t) d t.


\end{equation}





Как видно из формул (20), матрица $W$ размерности $n \times n$


зависит от неопределенных матриц $A$, $C$ и, по существу, неизвестна. Вектор


$g$ размерности $n \times 1$, напротив, определен однозначно измерениями


$y(t)$.





Полагая


\begin{equation}


W_0 = \int_0^{\theta} e^{t A_0'} C_0' C_0 e^{ t A_0} \, d t,


\end{equation}


по аналогии с [8] устанавливаем оценку


\begin{gather}


|W-W_0| \le \Delta W, \notag \\


\Delta W = \int_{0}^{\theta}| e^{ t A_0' } C_0'|


[( |C_0|+\Delta C )(e^{ t ( |A_0|+\Delta A ) } -


e^{ t |A_0|})+\Delta C | e^{ t A_0 } |] \, dt.


\end{gather} 


С учетом (22) систему линейных алгебраических уравнений (19), несколько


упрощая, можно считать интервальной. Следуя [8],


определим ее субуниверсальное решение $\wh x_0$ из системы уравнений


\begin{equation}


W_0 \wh x_0 = g


\end{equation}


с невырожденной [7] в силу (5) матрицей (21).





Используя соотношения (22), (23), найдем оценку невязки уравнения (19) на


субуниверсальном решении $ \wh x_0$. Имеем


\begin{equation}


|W \wh x_0 - g| = | W \wh x_0 - g -( W_0 \wh x_0 - g)| \le


|W-W_0|\,|\wh x_0| \le \Delta W |\wh x_0|


\end{equation}%24


для любых матриц $A$, $C$ из интервалов (3).





Очевидно, невязка (24) покоординатно уменьшается с уменьшением норм матриц


$\Delta A$, $\Delta C$ и в пределе при $\Delta A=0$, $\Delta C=0$


становится нулевой. В этом случае уравнение (23) дает точное значение вектора


$x_0$: $x_0=\wh x_0$.


Отметим попутно два необходимых для дальнейшего неравенства


\begin{gather}


|\wh x_0| \le L |x_0|, \\


%


L=|W_0^{-1}|\int_0^{\theta} | e^{t A_0'} C_0'|


( |C_0|+\Delta C )e^{ t (|A_0|+\Delta A)} \, d t;\notag\\


%


|x_0 - \wh x_0| \le |W_0^{-1}| \Delta W |x_0|,


\end{gather}


которые вытекают из соотношений (18), (20), (23) и (19), (22), (23)


соответственно.





Опишем вторую фазу активного управления системой (1) на отрезке времени


$[\theta, T]$, считая оценку $ \wh x_0$ вектора $x_0$ известной.


По аналогии с (8) полагаем


\begin{equation}


u(t)=B_0'e^{(T-t)A_0'} \wh z, \ \  \theta \le t \le T,


\end{equation} 


где $\wh z$ --- решение системы линейных алгебраических уравнений


\begin{equation}


\wh D_0 \wh z = \wh f_0


\end{equation}% 28


с матричными коэффициентами


\begin{equation}


\wh D_0=\int_0^{T-\theta} e^{\tau A_0}B_0 B_0'e^{\tau A_0'} \, d \tau,


\quad \wh f_0 = - e^{T A_0}\wh x_0


\end{equation}


размерности $n \times n$ и $n \times 1$ соответственно.





Оценим действие управления (27)--(29) на состояние $x(T)$


системы (1) с начальным условием $ x(\theta)= e^{\theta A}x_0 $


--- результатом свободного движения системы при $u=0$, $0 \le t \le


\theta$. По формуле Коши имеем


$$


x(T)=e^{(T-\theta) A} ( e^{\theta A}x_0 )+


\int_{\theta}^{T} e^{(T-t) A} B \wh u(t) \, dt.


$$


Отсюда, полагая


\begin{equation}


\wh D =\int_0^{T-\theta} e^{\tau A}B B_0'e^{\tau A_0'} \, d \tau,


\end{equation}


с учетом (27)--(29) и (11) получим


\begin{equation}


|x(T)|=|-f+\wh D \wh z |= | \wh D \wh z - f - (\wh D_0 \wh z -\wh f_0)|


\le |f-\wh f_0|+|\wh D -\wh D_0|\, |\wh z|.


\end{equation}





Оценим сверху каждое слагаемое в правой части неравенства (31). Для первого


из них на основании [8] и (26) имеем


\begin{gather}


|f-\wh f_0|=|e^{TA}x_0 - e^{TA_0}\wh x_0| =


|(e^{TA}-e^{TA_0})x_0+e^{TA_0}(x_0-\wh x_0)| \le \notag \\


%


\le |e^{TA}-e^{TA_0}|\,|x_0|+|e^{TA_0}|\, |x_0-\wh x_0| \le K |x_0|,\\


%


K= e^{T(|A_0|+\Delta A)}-e^{T|A_0|}+|e^{TA_0}|\,|W_0^{-1}| \Delta W.\notag


\end{gather}





Второе слагаемое оценим с использованием (29), (30), (28) и (25). В результате


получим


\begin{equation}


|\wh D -\wh D_0| \,|\wh z| \le


\Delta \wh D | \wh D_0^{-1} e^{TA_0}| L |x_0|,


\end{equation}


где матрица $\Delta \wh D$ получается из $\Delta D$ заменой $T$ на


$(T-\theta)$ в формуле (13):


$$


\Delta \wh D = \int_{0}^{T-\theta}


[(e^{ \tau ( |A_0|+\Delta A ) } -e^{\tau |A_0|})\,


( |B_0|+\Delta B )+| e^{ \tau A_0 } | \Delta B]


\, |B_0' e^{ \tau A_0' }| \, d\tau.


$$


Из (31) с использованием оценок (32), (33) находим окончательно


\begin{equation}


|x(T)| \le N |x_0|, \quad  N=K+\Delta \wh D |\wh D_0^{-1} e^{TA_0}| L.


\end{equation} %35





Как видно из соотношений (22), (25), (32), (34), определяющих матрицу


$N$, с уменьшением норм матриц $\Delta A$, $\Delta B$, $\Delta C$


норма матрицы $N$ стремится к нулю. В частности, при


$\Delta A=0$, $\Delta B=0$, $\Delta C=0$ получаем $N=0$, и из (34) следует


$x(T)=0$.


Следовательно, управление (27)--(29) переводит траекторию центральной системы


из начальной точки $x(\theta)=e^{\theta A_0}x_0$ в начало координат за время


$T-\theta$.





\section{Достаточное условие стабилизируемости}





Продолжим управление (27)--(29) на полуось $[0, \infty)$, полагая


\begin{gather}


\wh u(x_0,t)= 0, \ \ (k-1)T \le t < (k-1)T+\theta,\notag\\


\wh u(x_0,t)= B_0'e^{(kT-t)A_0'} \wh z^{k-1}, \ \ 


(k-1)T+\theta \le t < kT,\notag\\


\wh D_0 \wh z^{k-1} = \wh f_0^{k-1}, \ \


W_0 \wh x_0^{k-1} = g^{k-1}, \\


\wh f_0^{k-1} = - e^{TA_0} \wh x_0^{k-1},\notag \\


g^{k-1}= \int\limits_{(k-1)T}^{(k-1)T+\theta}


e^{(t-(k-1)T) A_0'} C_0' y(t) dt,\ \ k=1,2,\dotsc,\notag


\end{gather}


где $x(t)$, $y(t)$ --- решение системы (7). Так же, как при выводе неравенства


(34), легко устанавливается оценка


$$


|x(kT)| \le N |x((k-1)T)|, \ \ k=1,2,\dotsc


$$


Из нее в согласованных нормах получим


$$


\|x(kT) \| \le \|N\|^k \|x_0\|.


$$





Если $\|N\| < 1$, то для произвольных фиксированных $x_0$ из $R^n$ и


$A$, $B$, $C$ из интервалов (3) имеем $\|x(kT)\| \rightarrow 0$ при


$k \rightarrow \infty$. На отрезке времени $[(k-1)T, kT]$ справедливо


аналогичное (26) неравенство


$$


|x((k-1)T) - \wh x_0^{k-1}| \le |W_0^{-1}| \Delta W |x((k-1)T)|.


$$


Отсюда в силу сказанного заключаем $\wh x_0^{k-1} \rightarrow 0$ при


$k \rightarrow \infty$. В результате доказана


\medskip





{\bf Теорема.} {\it Пусть для системы} (1)--(3) {\it выполнены


предположения} (4), (5) {\it и матрица} $N$,


{\it определенная формулами} (22), (25), (32), (34), {\it удовлетворяет


условию} $\| N \| < 1$. {\it Тогда управление} (35) {\it обеспечивает


притяжение траекторий системы} (1)--(3) {\it к началу координат при любых


начальных состояниях} $x_0$ {\it из} $R^n$ {\it и матрицах} $A$, $B$, $C$


{\it из интервалов} (3).


\medskip





Продемонстрируем некоторые возможности теоремы


при упрощающих предположениях 


$$


C_0=E, \ \  \Delta C=0, \ \ \rank B_0=n


$$


и при любой неотрицательной матрице $\Delta A$.


Тогда условия $\|N\| <1$ теоремы


можно добиться уменьшением периода $T$ и нормы матрицы $\Delta B$.


Действительно, раскрывая формулу (34) для матрицы $N=N(T)$ и применив


поэлементно теорему о среднем для интеграла


$$


D_0 = \int_0^T e^{\tau A_0}B_0 B_0'e^{\tau A_0'} \, d \tau,


$$


в пределе получим 


$$


\lim\limits_{T \to 0} N(T)= \Delta B |B_0'|\, |(B_0 B_0')^{-1}|.


$$


Отсюда и вытекает требуемое.





Таким образом, в рассматриваемом случае за счет увеличения частоты коррекции


управления и уменьшения нормы матрицы $ \Delta B$ можно эффективно 


стабилизировать систему (1), (2) даже при больших интервалах неопределенности


матрицы $A$.





\section*{Заключение}





Отметим характерные особенности решения интервальной задачи стабилизации.





Метод упреждающего управления в сочетании с техникой интервального анализа


оказались достаточно эффективным средством анализа мультипликативных


неопределенностей в уравнениях движения и наблюдения управляемого объекта


и позволили получить детерминированный закон управления и достаточные


условия стабилизации. Упреждающее действие стабилизирующего управления


проявляется периодически на отрезках времени $[(k-1)T,kT]$, $k=1,2,\dotsc$


Управление ``работает'' на уменьшение норм векторов состояний $x(kT)$.


Техника интервального анализа дает возможность оценить ``недоработки''


управления по сравнению с действием на центральную систему.





Разделение фаз пассивного наблюдения и активного управления в процедуре


стабилизации имеет принципиальное значение для асимптотической устойчивости


положения равновесия. 


Наложение этих фаз приводит к появлению неоднородных членов


в уравнениях наблюдения (17) и, как следствие, к потере асимптотической


устойчивости. В этом случае, вообще говоря, можно гарантировать лишь


устойчивость положения равновесия.





Условия (4), (5) теоремы накладывают 


определенные требования на центры интервалов:


соответствующая центральная система должна быть управляемой и наблюдаемой.


Условие $\|N\| < 1$ можно интерпретировать двояко. С одной стороны,


при заданных $T$, $\theta$ оно выделяет в явном виде


``окрестность'' интервальных


систем (по отношению к центральной системе), которые поддаются стабилизации


приведенным способом. С другой стороны, это условие согласованности


показателя неопределенности системы


$( \| \Delta A \|+\| \Delta B \|+\| \Delta C \|)$


с необходимой частотой наблюдения и управления ($1/ \theta$ и $1/T$)


при стабилизации системы.





\begin{thebibliography}{99}





\bibitem{1}


Калман Р., Фалб П., Арбиб М. { \it Очерки по математической теории систем.}


-- М.: Мир, 1971. -- 400 с.


\bibitem{2}


Куржанский А.Б. {\it Управление и наблюдение в условиях неопределенности.}


-- М.: Наука, 1977. -- 390 с.


\bibitem{3}


Ащепков Л.Т., Стегостенко Ю.Б. { \it Стабилизация линейной дискретной


системы управления с интервальными коэффициентами} // Изв. вузов.


Математика. -- 1998. -- \No\,12. -- С.\,3--10.


\bibitem{4}


Ли Э., Маркус Л. {\it Основы теории оптимального управления.} -- М.: Наука,


1972. -- 572 с.


\bibitem{5}


Красовский Н.Н. { \it Теория управления движением. Линейные системы}.


-- М.: Наука, 1968. -- 476 с.


\bibitem{6}


Гантмахер Ф.Р. { \it Теория матриц.} -- М.: Наука, 1967. -- 576 с.


\bibitem{7}


Габасов Р.Ф., Кириллова Ф.М. { \it Оптимизация линейных систем.}


-- Минск: Изд-во Белорусск. гос. ун-та, 1973. -- 246 с.


\bibitem{8}


Ащепков Л.Т., Давыдов Д.В. { \it Стабилизация линейной стационарной


системы управления с интервальными коэффициентами} // Дальневосточный


матем. сб. -- Владивосток: Дальнаука, 1999. -- \No\,8. --  C.\,32--38.


\bibitem{9}


Понтрягин Л.С. { \it Обыкновенные дифференциальные уравнения.} --


М.: Наука, 1965. -- 332 с.


\bibitem{10}


Ащепков Л.Т. { \it Лекции по оптимальному управлению.} -- 


Владивосток: Изд-во Дальневосточн. ун-та, 1994. -- 308 с.





\end{thebibliography}





\vskip 1cm





\post{\it Институт прикладной математики}{\it Поступила}


\post{\it Дальневосточного отделения}{30.10.2000}


\post{\it Российскoй Академии наук}{}





\label{end}


\end{document}











Пример 2. Рассмотрим дифференциальные уравнения 


$$ \dot x_1 = x_2, \, \, \, \dot x_2 = -ax_1+bu, $$


описывающие движение механической системы "груз-пружина" (см., например, [10]).


Здесь $a=k/m>0, \, \, b=1/m>0; \, \, \, k$- коэффициент упругости пружины, 


$m$- масса груза, $ x_1$- расстояние от центра массы груза до начала координат, 


$ x_2$- скорость движения груза, $u$-действующая на груз сила.





Выясним с помощью теоремы следующий вопрос. Пусть фазовые состояния


$x_1(t), x_2(t)$ системы точно измеряются в дискретные моменты времени


$ t= 0, \, T, \, 2T, \, ... \, \, $. Каков должен быть интервал 


$ |a-a_0| \le \Delta a$ изменения параметра $a$ при фиксированном значении


параметра $b$, чтобы описанное ранее в двухэтапной процедуре управление


стабилизировало систему в окрестности положения равновесия


$ x_1=0, \, \, \, x_2=0$ 


при любом значении начального состояния 


$ x_1(0)=x_{10}, \, \, \, x_2(0)=x_{20}$. Другими словами, речь идет о 


пределах робастности стабилизирующего управления относительно 


коэффициента упругости пружины.





Полагая


$$


A_0= \left( 


\begin{array}{cc}


0 & 1 \\


-a_0 & 0 


\end{array}


\right), 


\, \, \, 


B_0= \left( 


\begin{array}{c}


0 \\


b 


\end{array}


\right), 


\, \, \, 


C_0=E, 


$$


$$


\Delta A= \left( 


\begin{array}{cc}


0 & 0 \\


\Delta a & 0 


\end{array}


\right), 


\, \, \, 


\Delta B= \left( 


\begin{array}{c}


0 \\


0 


\end{array}


\right), 


\, \, \, 


\Delta C=O,


$$


убеждаемся, что условия (4), (5) при $b>0$ выполнены. Матрица $N$ имеет


довольно громоздкое аналитическое выражение, содержащее тригонометрические


и гиперболические функции. Чтобы выписать простые достаточные условия


стабилизируемости, построим верхнюю поэлементную оценку $N_1$ матрицы $N$,


принимая следующие допущения: \\


$ 1) \sqrt{a_0+\Delta a} \le \sqrt{a_0}+\frac{\Delta a}{2 \sqrt{a_0}};$ \\


\\


$ 2) sh(x) \le ch(x) \le 1+x^2 \,\,\, \mbox{при} \, \, \, 0 \le x \le 2.9 ; $ \\


\\


$ 3) 1 < T \sqrt{a_0} \le 2.9$.





\newpage





В результате получим


$$


N_1= 5 \Delta a 


\left( 


\begin{array}{cc}


\frac{50T}{\sqrt{a_0}} & \frac{47T}{\sqrt{a_0}} \\


      &       \\


71T     & \frac{74T}{\sqrt{a_0}} 


\end{array}


\right).


\, \, \, 


$$





Примем за нормы матриц $N, \, \, \, N_1$ их максимальные по абсолютной величине 


элементы. Тогда неравенство


$ \| N \| \le \| N_1 \| < 1 $ 


будет гарантировано, если параметры 


$ T, \, \, \, a_0, \, \, \, \Delta a$ удовлетворяют системе неравенств


$$


71 T \Delta a \le 0.2 ; \, \, \, \, \, 74 \frac{T}{\sqrt{a_0}}\Delta a \le 0.2 ;


\, \, \, \, \, 1<T \sqrt{a_0} \le 2.9.


$$


Отсюда легко находим оценки


$$


\frac{1}{T^2} < a_0 \le \frac{8.41}{T^2}, \, \, \, \, \, 


\Delta a = \min \left\{ \frac{1}{355T}, \, \, \, \, \, 


\frac{\sqrt{a_0}}{370T} \right\}


$$


центра и полудлины искомого интервала


$ |a-a_0| \le \Delta a$. Заметим, что если в последних неравенствах


рассматривать $ a_0, \, \, \, \Delta a$ как заданные, то они


определяют допустимую частоту $1/T$ коррекции управления.


